\documentclass[10pt,journal,compsoc,onecolumn]{IEEEtran}
\usepackage[utf8]{inputenc}
\usepackage{amsmath,amsfonts,amssymb}
\usepackage{amsthm}
\usepackage{graphicx}
\usepackage{cite}
\usepackage{textcomp}
\usepackage{tikz}
\usetikzlibrary{arrows,positioning,shapes.geometric}
\usepackage{booktabs}
\usepackage{url}

\newtheorem{theorem}{Theorem}
\newtheorem{lemma}[theorem]{Lemma}
\newtheorem{definition}[theorem]{Definition}
\newtheorem{corollary}[theorem]{Corollary}

\begin{document}

\title{Graph Theory Foundations, Depth-First Search Mechanics, Structural Storage Representations, and Cycle Detection Algorithms in Directed and Undirected Graphs: A Formal Algorithmic Analysis}

\IEEEtitleabstractindextext{%
\begin{abstract}
This paper presents a formal, graduate-level mathematical formalization and algorithmic analysis of Graph Theory, Graph Storage Representations, Depth-First Search (DFS) Traversal Mechanics, and Cycle Detection Algorithms in both Undirected and Directed Graphs. We formalize graph properties, degree invariants, density metrics, and contrast Adjacency Matrix ($O(V^2)$ space) against Adjacency List ($O(V+E)$ space) representations with mathematical equivalence proofs and edge-list construction algorithms. We dissect Depth-First Search (DFS), demonstrating its recursive call stack mechanics and state space exploration across disconnected graph components. Furthermore, we address the challenge of cycle detection: for undirected graphs, we establish the Parent-Tracking invariant; for directed graphs, we formalize the distinction between reconvergent DAG paths (diamond structures) and true directed cycles, introducing the 3-Coloring State Machine (WHITE, GRAY, BLACK) and Back-Edge Theorem. Complete compilable Java reference implementations, execution trace tables, and a diagnostic taxonomy matrix are provided for all discussed algorithms.
\end{abstract}

\begin{IEEEkeywords}
binary trees, algorithms, data structures, computational complexity, tree traversal, recursive algorithms, formal analysis
\end{IEEEkeywords}
}

\maketitle
\IEEEdisplaynontitleabstractindextext
\IEEEpeerreviewmaketitle

\section{Introduction and Formal Graph Theory Specifications}
\label{sec:intro}

Graph structures constitute one of the most fundamental non-linear data structures in computer science and applied mathematics. Formally, a graph models pairwise relationships between discrete objects, serving as the core computational primitive for network routing, social network modeling, dependency resolution in software compilation, program analysis, and database query optimization.

Efficient processing of graph structures requires a clear mathematical taxonomy of graph types, invariant properties, and structural storage formats. In this section, we present a graduate-level formalization of graph theory primitives, topological classifications, vertex degree metrics, and density formulations.

\begin{definition}[Graph Topology]
A graph $G$ is a 2-tuple $G = (V, E)$ consisting of a finite, non-empty set of vertices (or nodes) $V$ and a set of edges $E$. The order of the graph is denoted by $|V| = n$, and the size of the graph is denoted by $|E| = m$.
\end{definition}

\begin{definition}[Directed vs. Undirected Graphs]
Graph topologies are broadly classified into two categories based on edge symmetry:
\begin{itemize}
    \item \textbf{Directed Graph (Digraph):} A directed graph $D = (V, E)$ consists of vertices $V$ and a set of ordered pairs $E \subseteq V \times V$. An edge $e = (u, v) \in E$ represents a unidirectional link from source vertex $u$ to target vertex $v$, where $(u, v) \neq (v, u)$.
    \item \textbf{Undirected Graph:} An undirected graph $G = (V, E)$ consists of vertices $V$ and a set of unordered pairs $E \subseteq \{\{u, v\} \mid u, v \in V, u \neq v\}$. An edge $e = \{u, v\}$ represents a bidirectional relation between $u$ and $v$, rendering $(u, v)$ and $(v, u)$ identical.
\end{itemize}
\end{definition}

\begin{definition}[Weighted vs. Unweighted Graphs]
Edge sets may be augmented with a numeric weight mapping:
\begin{equation}
w: E \to \mathbb{R}
\end{equation}
In an unweighted graph, all edge weights are implicitly uniform, such that $w(e) = 1$ for all $e \in E$. In a weighted graph, $w(e)$ represents domain-specific metrics such as distance, latency, capacity, or cost.
\end{definition}

\begin{definition}[Graph Connectedness and Components]
The connectivity structure of a graph dictates reachable paths across vertex pairs:
\begin{itemize}
    \item \textbf{Undirected Connectivity:} An undirected graph is \emph{connected} if for every pair of vertices $u, v \in V$, there exists a path $P = (u = v_0, v_1, v_2, \dots, v_k = v)$ such that $\{v_i, v_{i+1}\} \in E$ for all $0 \le i < k$. A disconnected undirected graph partitions into $k \ge 2$ maximal connected subgraphs termed \emph{Connected Components} $C_1, C_2, \dots, C_k$.
    \item \textbf{Directed Connectivity:} A directed graph is \emph{strongly connected} if there exists a directed path from $u$ to $v$ AND from $v$ to $u$ for every pair $u, v \in V$. It is \emph{weakly connected} if replacing all directed edges with undirected edges yields a connected undirected graph.
\end{itemize}
\end{definition}

\begin{definition}[Vertex Degree and In-Out Invariants]
The degree of a vertex measures its local edge connectivity:
\begin{itemize}
    \item For an undirected graph $G=(V, E)$, the degree of a vertex $v \in V$, denoted $\deg(v)$, is the number of edges incident to $v$:
    \begin{equation}
    \deg(v) = |\{ u \in V \mid \{v, u\} \in E \}|
    \end{equation}
    \item For a directed graph $D=(V, E)$, the degree partitions into \emph{in-degree} $\deg^-(v)$ and \emph{out-degree} $\deg^+(v)$:
    \begin{equation}
    \deg^-(v) = |\{ u \in V \mid (u, v) \in E \}|, \quad \deg^+(v) = |\{ u \in V \mid (v, u) \in E \}|
    \end{equation}
    The total degree is given by $\deg(v) = \deg^-(v) + \deg^+(v)$.
\end{itemize}
\end{definition}

\begin{theorem}[Handshaking Lemma]
\label{thm:handshaking}
For any finite graph, the sum of vertex degrees is directly proportional to the total edge count:
\begin{enumerate}
    \item For an undirected graph $G = (V, E)$:
    \begin{equation}
    \sum_{v \in V} \deg(v) = 2|E|
    \end{equation}
    \item For a directed graph $D = (V, E)$:
    \begin{equation}
    \sum_{v \in V} \deg^-(v) = \sum_{v \in V} \deg^+(v) = |E|
    \end{equation}
\end{enumerate}
\end{theorem}

\begin{proof}
Every undirected edge $e = \{u, v\}$ contributes exactly 1 to $\deg(u)$ and 1 to $\deg(v)$, thereby contributing 2 to the sum $\sum_{v \in V} \deg(v)$. Summing across all $|E|$ edges yields $2|E|$. Similarly, each directed edge $e = (u, v)$ contributes exactly 1 to $\deg^+(u)$ and 1 to $\deg^-(v)$. Summing across all directed edges yields $\sum \deg^+(v) = |E|$ and $\sum \deg^-(v) = |E|$.
\end{proof}

\begin{definition}[Simple Graphs, Self-Loops, and Density]
A \emph{simple graph} contains no self-loops ($(v, v) \notin E$) and no parallel (duplicate) edges between any node pair. Graph density $\rho$ expresses the ratio of existing edges to maximum possible edges:
\begin{equation}
\rho_{\text{undirected}} = \frac{2|E|}{|V|(|V|-1)}, \quad \rho_{\text{directed}} = \frac{|E|}{|V|(|V|-1)}
\end{equation}
A graph is \emph{sparse} if $|E| = O(|V|)$ (where $\rho \to 0$ as $|V| \to \infty$) and \emph{dense} if $|E| = \Theta(|V|^2)$ (where $\rho = \Theta(1)$).
\end{definition}

% -----------------------------------------------------------------------------
\section{Graph Storage Representations and Structural Trade-offs}
\label{sec:representations}

Choosing an appropriate memory representation for a graph is a critical engineering decision that dictates spatial overhead and computational complexity for graph operations. The two primary structural formats are the \textbf{Adjacency Matrix} and the \textbf{Adjacency List}.

\subsection{Adjacency Matrix Representation}

An Adjacency Matrix represents a graph $G = (V, E)$ with $|V| = N$ as a two-dimensional square matrix $M$ of size $N \times N$, indexed from $0$ to $N-1$ (or $1$ to $N$).

\begin{equation}
M[u][v] = \begin{cases}
1 \text{ (or } w(u, v) \text{)}, & \text{if } (u, v) \in E \\
0 \text{ (or } \infty \text{)}, & \text{if } (u, v) \notin E
\end{cases}
\end{equation}

For undirected graphs, $M$ is symmetric across the main diagonal ($M[u][v] = M[v][u]$). For directed graphs, $M[u][v] = 1$ indicates an edge originating at $u$ and terminating at $v$.

\textbf{Properties of Adjacency Matrix:}
\begin{itemize}
    \item \textbf{Space Complexity:} $\Theta(V^2)$ memory words, independent of $|E|$.
    \item \textbf{Edge Existence Query $(u, v) \in E$:} $O(1)$ constant time lookup via $M[u][v]$.
    \item \textbf{Neighbor Exploration:} Iterating over all out-neighbors of $u$ requires scanning row $M[u]$, taking $\Theta(V)$ time regardless of $\deg^+(u)$.
\end{itemize}

\subsection{Adjacency List Representation}

An Adjacency List represents a graph $G = (V, E)$ as an array of dynamic headers of size $N$, where entry $L[u]$ contains a dynamic collection (linked list or dynamic array) of all adjacent target vertices $v$ such that $(u, v) \in E$.

\textbf{Properties of Adjacency List:}
\begin{itemize}
    \item \textbf{Space Complexity:} $\Theta(V + E)$ memory words ($\Theta(V + 2E)$ for undirected, $\Theta(V + E)$ for directed).
    \item \textbf{Edge Existence Query $(u, v) \in E$:} $O(\deg^+(u))$ time, requiring a linear scan of list $L[u]$.
    \item \textbf{Neighbor Exploration:} Iterating over out-neighbors of $u$ takes $O(\deg^+(u))$ time, which is strictly optimal.
\end{itemize}

\subsection{Comparative Analysis and Trade-off Matrix}

Table~\ref{tab:rep_comparison} summarizes the asymptotic trade-offs between Adjacency Matrix and Adjacency List representations.

\begin{table}[htbp]
\caption{Asymptotic Comparison of Graph Storage Representations}
\label{tab:rep_comparison}
\centering
\begin{tabular}{lcc}
\toprule
\textbf{Operation / Property} & \textbf{Adjacency Matrix} & \textbf{Adjacency List} \\
\midrule
Space Complexity & $\Theta(V^2)$ & $\Theta(V + E)$ \\
Edge Existence Query $(u, v) \in E$ & $\Theta(1)$ & $O(\deg(u))$ \\
Enumerate Out-Neighbors of $u$ & $\Theta(V)$ & $\Theta(\deg(u))$ \\
Enumerate In-Neighbors of $v$ & $\Theta(V)$ & $\Theta(V + E)$ (or $\Theta(\deg^-(v))$ w/ transpose) \\
Add Edge $(u, v)$ & $\Theta(1)$ & $\Theta(1)$ \\
Delete Edge $(u, v)$ & $\Theta(1)$ & $O(\deg(u))$ \\
Optimal Application Domain & Dense Graphs ($E \approx V^2$) & Sparse Graphs ($E \ll V^2$) and General DSA \\
\bottomrule
\end{tabular}
\end{table}

\subsection{Formal Java Implementation: Edge List Construction}

\begin{verbatim}
import java.util.*;

public class GraphRepresentation {
    // Construct Adjacency Matrix for Undirected Graph
    public static int[][] buildAdjMatrixUndirected(int n, int[][] edges) {
        int[][] matrix = new int[n + 1][n + 1];
        for (int[] edge : edges) {
            int u = edge[0];
            int v = edge[1];
            matrix[u][v] = 1;
            matrix[v][u] = 1;
        }
        return matrix;
    }

    // Construct Adjacency List for Directed / Undirected Graph
    public static List<List<Integer>> buildAdjList(int n, int[][] edges, boolean isDirected) {
        List<List<Integer>> adj = new ArrayList<>();
        for (int i = 0; i <= n; i++) {
            adj.add(new ArrayList<>());
        }
        for (int[] edge : edges) {
            int u = edge[0];
            int v = edge[1];
            adj.get(u).add(v);
            if (!isDirected) {
                adj.get(v).add(u);
            }
        }
        return adj;
    }
}
\end{verbatim}

% -----------------------------------------------------------------------------
\section{Depth-First Search (DFS) Traversal Mechanics and Call Stack Analysis}
\label{sec:dfs}

Depth-First Search (DFS) is a fundamental traversal paradigm for systematic state-space exploration of graph structures. While Breadth-First Search (BFS) explores vertices level-by-level using a Queue, DFS dives as deep as possible along each branch before backtracking using a Stack (either implicit via recursive call frames or explicit via data structures).

\subsection{Algorithmic Invariants and Visited Tracking}

Traversing a graph with cycles introduces the risk of infinite recursion loops. To guarantee termination and linear time bounds, DFS maintains a boolean state array \texttt{visited[]} of size $N+1$.

\begin{definition}[DFS Traversal Invariant]
During DFS execution on graph $G = (V, E)$:
\begin{enumerate}
    \item A vertex $u \in V$ is marked \texttt{visited[u] = true} at the exact moment its recursive frame is initiated.
    \item A vertex $u$ is processed at most once across the entire execution.
    \item The outer traversal loop iterates over all $v \in V$, triggering a new DFS tree root whenever \texttt{visited[v] == false}. This guarantees complete traversal across disconnected components (forming a \emph{DFS Forest}).
\end{enumerate}
\end{definition}

\subsection{Formal Complexity Analysis}

\begin{theorem}[DFS Time and Space Complexity]
\label{thm:dfs_complexity}
For a graph $G = (V, E)$ represented via Adjacency List, Depth-First Search completes in:
\begin{equation}
T(V, E) = \Theta(V + E) \text{ time}, \quad S(V) = O(V) \text{ auxiliary space}
\end{equation}
\end{theorem}

\begin{proof}
We prove time complexity via Aggregate Analysis:
\begin{enumerate}
    \item The outer forest loop executes $V$ times, checking \texttt{visited[i]}. This contributes $\Theta(V)$ time.
    \item Each vertex $u \in V$ is marked \texttt{visited[u] = true} exactly once. Thus, \texttt{dfs(u)} is called at most once per vertex across the lifetime of the program.
    \item Inside \texttt{dfs(u)}, the neighbor traversal loop runs $\deg(u)$ times ($\deg^+(u)$ for directed). Summing over all vertices:
    \begin{equation}
    \sum_{u \in V} O(1 + \deg(u)) = O(V) + O\left(\sum_{u \in V} \deg(u)\right) = O(V + 2E) = O(V + E)
    \end{equation}
    Hence, total execution time is $\Theta(V + E)$.
    \item Auxiliary space is bounded by the maximum recursion depth (call stack frames) plus the \texttt{visited[]} array of size $V + 1$. In the worst-case (a linear path graph $1 \to 2 \to \dots \to N$), call stack depth is $V$. Thus, $S(V) = O(V)$.
\end{enumerate}
\end{proof}

\subsection{Java Implementation for Graph Traversal}

\begin{verbatim}
import java.util.*;

public class GraphDFS {
    public static void dfs(int currNode, List<List<Integer>> adj, boolean[] visited, List<Integer> order) {
        visited[currNode] = true;
        order.add(currNode);

        for (int neighbor : adj.get(currNode)) {
            if (!visited[neighbor]) {
                dfs(neighbor, adj, visited, order);
            }
        }
    }

    public static List<Integer> traverseGraph(int n, List<List<Integer>> adj) {
        boolean[] visited = new boolean[n + 1];
        List<Integer> traversalOrder = new ArrayList<>();

        // Forest loop for disconnected components
        for (int i = 1; i <= n; i++) {
            if (!visited[i]) {
                dfs(i, adj, visited, traversalOrder);
            }
        }
        return traversalOrder;
    }
}
\end{verbatim}

% -----------------------------------------------------------------------------
\section{Formal Analysis of Cycle Detection in Undirected Graphs}
\label{sec:undirected_cycle}

Detecting cycles is a foundational graph problem required to verify tree structures, detect deadlocks in resource allocation, and validate network topologies.

\begin{definition}[Undirected Cycle]
A cycle in an undirected graph $G = (V, E)$ is a sequence of vertices $(v_0, v_1, v_2, \dots, v_k = v_0)$ such that $k \ge 3$, $\{v_i, v_{i+1}\} \in E$ for $0 \le i < k$, and all intermediate vertices $v_0, v_1, \dots, v_{k-1}$ are distinct.
\end{definition}

\subsection{The Parent Tracking Strategy}

In an undirected graph, every edge $\{u, v\}$ is symmetric, appearing in both $L[u]$ (pointing to $v$) and $L[v]$ (pointing to $u$). If DFS moves from $u \to v$, when exploring neighbors of $v$, it immediately sees $u$. If $u$ is treated as an ordinary visited neighbor, DFS would falsely classify the immediate reverse path $v \to u$ as a cycle.

To eliminate this trivial false positive, DFS must pass a \texttt{parent} node parameter along each recursive frame.

\begin{definition}[Undirected Cycle Detection Invariant]
During DFS on an undirected graph, if a vertex $u$ encounters an adjacent neighbor $v$ such that:
\begin{equation}
\text{visited}[v] == \text{true} \quad \text{AND} \quad v \neq \text{parent}
\end{equation}
then a non-trivial cycle exists in the graph.
\end{definition}

\subsection{Lemma and Correctness Proof}

\begin{lemma}[Undirected Cycle Characterization]
\label{lem:undirected_cycle}
An undirected graph $G = (V, E)$ contains a cycle if and only if DFS encounters an already-visited vertex $v$ that is distinct from the immediate parent of the current vertex $u$.
\end{lemma}

\begin{proof}
$(\Leftarrow)$ Suppose DFS at vertex $u$ encounters neighbor $v$ with \texttt{visited[v] == true} and $v \neq \text{parent}$. Since $v$ is already visited, $v$ is an ancestor of $u$ in the current DFS tree. Therefore, there exists a path $P_{v \to \dots \to u}$ in the DFS tree from $v$ to $u$. Concatenating edge $\{u, v\}$ to $P_{v \to \dots \to u}$ forms a closed simple path $v \to \dots \to u \to v$ of length $\ge 3$, establishing a cycle.

$(\Rightarrow)$ Conversely, suppose $G$ contains a cycle $C = (v_1, v_2, \dots, v_k, v_1)$. Let $v_i$ be the first vertex in $C$ discovered by DFS. DFS will traverse the cycle edges. Since $k \ge 3$, when DFS reaches the last vertex $v_j$ in the cycle adjacent to $v_i$, $v_i$ will be already marked visited, and $v_i$ will not be the immediate parent of $v_j$. Thus, DFS will detect the cycle condition.
\end{proof}

\subsection{Java Implementation: Undirected Cycle Detection}

\begin{verbatim}
import java.util.*;

public class UndirectedCycleDetection {
    private static boolean dfsCheckCycle(int node, int parent, List<List<Integer>> adj, boolean[] visited) {
        visited[node] = true;

        for (int neighbor : adj.get(node)) {
            if (!visited[neighbor]) {
                if (dfsCheckCycle(neighbor, node, adj, visited)) {
                    return true;
                }
            } else if (neighbor != parent) {
                // Visited neighbor that is NOT parent -> Cycle detected!
                return true;
            }
        }
        return false;
    }

    public static boolean hasCycle(int n, List<List<Integer>> adj) {
        boolean[] visited = new boolean[n + 1];

        // Forest loop to handle disconnected graphs
        for (int i = 1; i <= n; i++) {
            if (!visited[i]) {
                if (dfsCheckCycle(i, -1, adj, visited)) {
                    return true;
                }
            }
        }
        return false;
    }
}
\end{verbatim}

% -----------------------------------------------------------------------------
\section{Formal Analysis of Cycle Detection in Directed Graphs}
\label{sec:directed_cycle}

Cycle detection in directed graphs requires a significantly more sophisticated state model than undirected cycle detection.

\subsection{The Reconvergent Path Fallacy vs. True Directed Cycles}

A fundamental pitfall when extending undirected cycle detection to directed graphs is confusing \emph{reconvergent paths in Directed Acyclic Graphs (DAGs)} with \emph{true directed cycles}.

Starting DFS at $A$ in a diamond DAG ($A \to B \to D$ and $A \to C \to D$): DFS explores path $A \to B \to D$. Vertex $D$ is marked \texttt{visited[D] = true}. DFS backtracks to $A$. DFS then explores path $A \to C \to D$. When at $C$, neighbor $D$ is inspected. \texttt{visited[D]} is \texttt{true}.

If we applied the single \texttt{visited[]} array logic here, $D$ being visited would trigger a false cycle positive. However, there is NO directed cycle! The paths $A \to B \to D$ and $A \to C \to D$ merely \emph{reconverge} at $D$. $D$ is no longer on the current execution stack when $C$ visits it.

\subsection{State Tracking: 3-Coloring State Machine}

To correctly identify directed cycles, we must distinguish between:
\begin{enumerate}
    \item Vertices that are currently active on the recursive execution stack (\emph{in-path}).
    \item Vertices that have been completely explored and popped from the call stack.
\end{enumerate}

This leads to the formal \textbf{3-Color State Machine}:

\begin{definition}[3-Coloring States]
Each vertex $u \in V$ resides in exactly one of three states during DFS:
\begin{itemize}
    \item \textbf{State 0 / WHITE / Unvisited:} Vertex $u$ has not yet been discovered by DFS.
    \item \textbf{State 1 / GRAY / Visiting (\texttt{pathVisited[u] = true}):} Vertex $u$ has been discovered, and its recursive stack frame \texttt{dfs(u)} is currently ACTIVE on the call stack.
    \item \textbf{State 2 / BLACK / Visited (\texttt{visited[u] = true} and \texttt{pathVisited[u] = false}):} Vertex $u$ and all its downstream descendants have been fully explored. \texttt{dfs(u)} has completed and popped from the call stack.
\end{itemize}
\end{definition}

\subsection{Back-Edge Theorem and Correctness Proof}

\begin{theorem}[Back-Edge Characterization of Directed Cycles]
\label{thm:back_edge}
A directed graph $D = (V, E)$ contains a directed cycle if and only if a Depth-First Search traversal encounters a \textbf{back-edge} --- an edge $(u, v) \in E$ targeting a vertex $v$ that is currently in State 1 / GRAY (\texttt{pathVisited[v] == true}).
\end{theorem}

\begin{proof}
$(\Leftarrow)$ Suppose DFS at vertex $u$ encounters edge $(u, v)$ with $v$ in State 1 / GRAY. Since $v$ is in State 1, $v$ is currently active on the execution stack. Therefore, $v$ is an ancestor of $u$ in the current DFS call stack, meaning there exists a directed path $P_{v \to \dots \to u}$ on the call stack. Adding edge $(u, v)$ closes the loop $v \to \dots \to u \to v$, forming a valid directed cycle.

$(\Rightarrow)$ Suppose $D$ contains a directed cycle $C = (v_0, v_1, \dots, v_k = v_0)$. Let $v_i$ be the first vertex of $C$ discovered by DFS. At the moment $v_i$ is visited, it enters State 1 / GRAY. By DFS mechanics, all vertices reachable from $v_i$ along unvisited paths will be explored before \texttt{dfs(v\_i)} finishes. Since $C$ is a directed cycle, $v_{i-1}$ (the predecessor of $v_i$ in $C$) is reachable from $v_i$. When DFS reaches $v_{i-1}$, \texttt{dfs(v\_i)} is still active on the stack (State 1 / GRAY). Thus, inspecting edge $(v_{i-1}, v_i)$ yields a target in State 1, discovering a back-edge.
\end{proof}

\subsection{Production Java Implementations}

We provide two formal implementations: (1) using dual boolean arrays (\texttt{visited[]} + \texttt{pathVisited[]}), and (2) using a single integer state array (\texttt{state[]} with values 0, 1, 2).

\begin{verbatim}
import java.util.*;

public class DirectedCycleDetection {

    // Approach 1: Dual Boolean Arrays (visited[] and pathVisited[])
    private static boolean dfsDualArray(int u, List<List<Integer>> adj, 
                                        boolean[] visited, boolean[] pathVisited) {
        visited[u] = true;
        pathVisited[u] = true; // Push to recursion stack (GRAY)

        for (int v : adj.get(u)) {
            if (!visited[v]) {
                if (dfsDualArray(v, adj, visited, pathVisited)) {
                    return true;
                }
            } else if (pathVisited[v]) {
                // Back-edge detected! v is currently on recursion stack
                return true;
            }
        }

        pathVisited[u] = false; // Backtrack / Pop from stack (BLACK)
        return false;
    }

    // Approach 2: 3-Color State Array (0 = UNVISITED, 1 = VISITING, 2 = VISITED)
    private static final int UNVISITED = 0;
    private static final int VISITING = 1;
    private static final int VISITED = 2;

    private static boolean dfsThreeColor(int u, List<List<Integer>> adj, int[] state) {
        state[u] = VISITING; // Mark GRAY

        for (int v : adj.get(u)) {
            if (state[v] == VISITING) {
                // Target v is currently being visited -> Directed Cycle!
                return true;
            }
            if (state[v] == UNVISITED) {
                if (dfsThreeColor(v, adj, state)) {
                    return true;
                }
            }
        }

        state[u] = VISITED; // Mark BLACK
        return false;
    }

    public static boolean hasDirectedCycle(int n, List<List<Integer>> adj) {
        int[] state = new int[n + 1];

        for (int i = 1; i <= n; i++) {
            if (state[i] == UNVISITED) {
                if (dfsThreeColor(i, adj, state)) {
                    return true;
                }
            }
        }
        return false;
    }
}
\end{verbatim}

% -----------------------------------------------------------------------------
\section{Comparative Synthesis, Edge Cases, and Diagnostic Taxonomy Matrix}
\label{sec:synthesis}

To synthesize our findings, we present a unified taxonomy matrix in Table~\ref{tab:taxonomy_matrix} contrasting cycle detection paradigms across graph structural variations.

\begin{table}[htbp]
\caption{Diagnostic Taxonomy Matrix for Graph Traversal and Cycle Detection}
\label{tab:taxonomy_matrix}
\centering
\begin{tabular}{p{2.2cm}p{3.2cm}p{3.2cm}p{2.2cm}p{2.5cm}}
\toprule
\textbf{Graph Domain} & \textbf{Cycle Invariant Condition} & \textbf{State Data Structure} & \textbf{Time Bound} & \textbf{Primary Edge Cases} \\
\midrule
Undirected Graph & Visited neighbor $v \neq \text{parent}$ & \texttt{visited[]} array + \texttt{parent} variable & $O(V + E)$ & Parallel edges, self-loops, 2-node cycles \\
\midrule
Directed Graph & Edge $(u, v)$ targeting $v$ with State == GRAY & Dual \texttt{visited[]} + \texttt{pathVisited[]} OR 3-Color \texttt{state[]} & $O(V + E)$ & Reconvergent DAG paths (Diamonds), self-loops \\
\midrule
Directed Acyclic Graph (DAG) & Guaranteed NO Back-Edges ($\forall$ edges target WHITE or BLACK) & Topological Sort / Kahn's In-Degree Queue & $O(V + E)$ & Disconnected components, isolated nodes \\
\bottomrule
\end{tabular}
\end{table}

\subsection{Analysis of Edge Cases and Boundary Conditions}

\begin{enumerate}
    \item \textbf{Self-Loops ($(u, u) \in E$):} A self-loop represents a cycle of length 1. In directed graphs, \texttt{dfs(u)} inspects neighbor $u$, sees state GRAY, and immediately returns \texttt{true}. In undirected graphs, $u$ is visited and $u \neq \text{parent}$ (since initial parent is -1 or caller), correctly returning \texttt{true}.
    \item \textbf{Disconnected Graphs:} Handled seamlessly by the outer forest loop ($\forall i \in [1, N]$), ensuring that unreachable components are independently evaluated.
    \item \textbf{Parallel Edges in Undirected Graphs:} If two edges exist between $u$ and $v$, moving $u \to v$ via edge 1 and back $v \to u$ via edge 2 constitutes a valid cycle of length 2. In simple graph adjacency lists, duplicate entries must be retained or handled via edge ID tracking.
\end{enumerate}

\section{Conclusion}
\label{sec:conclusion}

This paper has presented a rigorous mathematical formalization and computational analysis of Graph Storage Representations, Depth-First Search mechanics, and Cycle Detection paradigms. We established the $O(V^2)$ vs $O(V+E)$ spatial trade-offs between Adjacency Matrices and Adjacency Lists, analyzed DFS activation stack frame transitions, proved the Undirected Parent-Tracking Invariant, and resolved the Directed Reconvergent Path Fallacy via the 3-Coloring State Machine and Back-Edge Theorem. The presented Java reference implementations and diagnostic taxonomy matrix provide a robust foundation for theoretical graph analysis and high-performance algorithmic execution.


\end{document}
